\section{2024-2025学年线性代数II（H）期中}

\begin{center}
    任课老师：吴志祥\hspace{4em} 考试时长：120分钟

    考试时间：2025年4月23日（星期三）10:00-12:00
\end{center}

\begin{enumerate}
    \item （10 分）求过点 $(1, 0, 2)$ 与平面 $3x - y + 2z + 2 = 0$ 平行，且与直线 $\dfrac{x-1}{4}=\dfrac{y+1}{-2}=\dfrac{z}{1}$ 相交的直线方程.
    \item （10 分）设 $V$ 为有限维线性空间，$T$ 是 $V$ 到 $V$ 的线性映射。若 $T^2 - 3T + 2I = 0$，证明：
    \[V = \operatorname{null}(T - I) \oplus \operatorname{range}(T - I).\]
    \item （10 分）设 $v_1, \ldots, v_n$ 为有限维线性空间 $V$ 的一组基，$\varphi_1, \ldots, \varphi_n$ 为其对偶基. 设 $\psi \in V^*$，证明：
    \[\psi = \psi(v_1)\varphi_1 + \cdots + \psi(v_n)\varphi_n.\]
    \item （10 分）设 $p \in \mathbf{C}[x]$，$p \neq 0$，令 $U = \{pq \mid q \in \mathcal{P}(\mathbf{C})\}$.
    \begin{enumerate}
        \item 证明：$\dim(\mathcal{P}(\mathbf{C}) / U) = \deg p$.
        \item 求 $\mathcal{P}(\mathbf{C}) / U$ 的一组基.
    \end{enumerate}
    \item （10 分）设 $V$ 为有限维内积空间，$S, T \in L(V)$. 称 $S,T$ 可同时对角化，若存在 $V$ 的一组基，使得 $S$ 和 $T$ 在这组基下的矩阵都是对角矩阵. 若 $S$ 和 $T$ 可对角化，证明：它们可同时对角化当且仅当 $ST = TS$.
    \item （10 分）设 $A$ 为 $n$ 阶实方阵，记 $A^\mathrm{T}$ 为 $A$ 的转置。 证明 $\dim E(\lambda, A) = \dim E(\lambda, A^\mathrm{T})$.
    \item （10 分）设 $W_1, W_2$ 为有限维内积空间的子空间。证明：
    \[(W_1 + W_2)^\perp = W_1^\perp \cap W_2^\perp;\quad (W_1 \cap W_2)^\perp = W_1^\perp + W_2^\perp.\]
    \item （10 分）考虑 $\mathbf{R}^n$ 上的标准内积，设列向量 $x_1,\ldots,x_p\in\mathbf{R}^n$，$\displaystyle\sum_{i=1}^p x_i=0$. 记 $X=(x_1,\ldots,x_p)$，设 $XX^\mathrm{T}$ 有 $n$ 个不同的正特征值，给定 $k\in\mathbf{N}_+$，$k<n$，考虑 PCA 算法：
    \begin{enumerate}
        \item $C=XX^\mathrm{T}$;
        \item 取 $\lambda_1\geqslant\cdots\geqslant\lambda_k$ 为 $C$ 最大的 $k$ 个特征值，$v_1,\ldots,v_k$ 为对应的单位特征向量;
        \item 令 $V=\spa(v_1,\ldots,v_k)$，$U=(v_1,\ldots,v_k)$，$P_V=UU^\mathrm{T}$;
        \item 令 $y_i=P_V x_i$，$i=1,\ldots,n$.
    \end{enumerate}
    证明：
    \begin{enumerate}
        \item $\{v_1,\ldots,v_k\}$ 为 $V$ 的规范正交基，且 $P_V$ 为 $\mathbf{R}^n$ 到 $V$ 的正交投影.
        \item PCA 算法最小化投影误差. 即 $\forall W$ 为 $\mathbf{R}^n$ 的 $k$ 维子空间，$P_W$ 为正交投影， $\displaystyle\sum_{i=1}^p \|x_i-y_i\|^2 \leqslant \displaystyle\sum_{i=1}^p \|x_i-P_W x_i\|^2$.
    \end{enumerate}
    \item （20 分）试判断下列命题的真伪。若命题为真，请给出简要证明；若命题为假，请举出反例。
    \begin{enumerate}
        \item 任意空间向量 $\vec{a}, \vec{b}, \vec{c}$ 满足 $(\vec{a}\times\vec{b})\times\vec{c} = \vec{a}\times(\vec{b}\times\vec{c})$.
        \item 每个非常数的复系数多项式都有零点.
        \item 设 $T \in L(V)$，存在 $V$ 的一组基 $\{v_1, \ldots, v_n\}$，使得 $T$ 关于这组基的矩阵是对角矩阵.
        \item 设 $V = C([0,1])$（即 $V$ 为 $[0,1]$ 区间上连续函数全体构成的线性空间），$\langle f, g \rangle = \displaystyle\int_0^{\frac{1}{2}} f(t)g(t)\,\mathrm{d}t$ 为 $V$ 上的内积.
    \end{enumerate}
\end{enumerate}

\clearpage
